% =====================================================================
\section{Dopamine and psychosis}
\fig{fig_dopamine.png}{The path from the hippocampus to the dopamine cells, after the work of Grace's group. An arrow is excitation, a bar at the end inhibition.}
Damage to PV cells does not by itself explain why voices and delusions arise. The bridge to the symptoms is taken to be \emph{dopamine}, the transmitter with which the brain marks events as significant. If its release is excessive and disorderly, random events appear significant, and delusion arises as an attempt to explain them.
\facts
\begin{itemize}
\item All antipsychotics in use block dopamine receptors of one type (D2). The effect appears when about 65\,\% of the receptors are occupied; above 78\,\% motor side effects begin.\tP
\item In a synthesis of 44 studies, dopamine production and release are elevated in patients, while the number of receptors is almost unchanged.\tP
\item Substances that enhance dopamine release cause or exacerbate psychosis.\tU
\item Grace's scheme: weakening of PV cells in the hippocampus leaves it over-excited; through two inhibitory links this increases the fraction of dopamine cells ready to respond to a signal (fig.~9).\tS
\end{itemize}
\begin{ineq}
Damage to inhibitory cells changes the gain of the dopamine system. The project's model does not have this layer: it ends at the state of a single cell. Symptoms are a separate future problem, with a time scale of its own, of seconds.
\end{ineq}
